\sect{Существующие применения RL к отслеживанию партитуры и
аккомпанементу}{rl-music}
%
Известные авторам работы, в которых обучение с подкреплением
применяется непосредственно к музыкальному аккомпанементу или к
отслеживанию партитуры, ниже разобраны последовательно. Для
каждой работы указаны постановка MDP (состояние, действие,
функция вознаграждения), архитектура агента, алгоритм обучения и
характер новизны. Цель обзора~--- не пересказать содержание
оригинальных статей, а сделать возможным точное и честное
позиционирование теоретического предложения~\Secref{proposal}.

\paragraph*{Дорфер, Хенкель и Видмер 2018: отслеживание
партитуры как RL-игра.} Принципиально первая работа этого
направления~\cite{Dorfer2018RL} формализует задачу как
\textit{мультимодальный} MDP. Состояние~$s_t$ объединяет два
сенсорных канала: окно изображения нотного стана размера
$80\times256$ пикселей и логарифмическую спектрограмму последних
$\approx 2$\,секунд аудио ($78\times 40$). Дополнительно для
обеспечения марковости в состояние включаются \textit{приращения
относительно предыдущего шага}~--- это стандартный приём,
позволяющий модели опираться на локальную динамику. Действие
$a_t\in\{-\Delta v_{\text{pxl}}, 0, +\Delta v_{\text{pxl}}\}$~---
дискретное изменение скорости продвижения курсора по
изображению партитуры. Архитектура агента изображена на
\Figref{pipeline}~--- в Дорфер-постановке RL заменяет весь
трекер целиком, что и отличает её от нашего предложения в
\Secref{proposal}.

\par\addvspace{2pt}
{\centering
\refstepcounter{figure}\figlab{pipeline}%
\resizebox{\columnwidth}{!}{% =====================================================================
% TikZ figure: end-to-end real-time pipeline (BLACK & WHITE only)
% Compact 3+3 layout to fit a single REVTeX two-column column.
% Used in: sections/05-realtime-pipeline.tex
% =====================================================================
\begin{tikzpicture}[
    node distance = 4mm and 4mm,
    every node/.append style={font=\scriptsize},
    block/.style ={rectangle, draw=black, thick, rounded corners=2pt,
                   minimum height=8mm, minimum width=18mm,
                   align=center, font=\scriptsize},
    arrow/.style ={-Stealth, thick}]

  \node[block]                 (in)    {MIDI / audio\\input};
  \node[block, right=of in]    (clean) {Cleaning\\\& buffer};
  \node[block, right=of clean] (hmm)   {HMM\\tracker};

  \node[block, below=8mm of hmm]   (oltw) {OLTW\\fall-back};
  \node[block, left=of oltw]       (rl)   {RL tempo\\anticip.};
  \node[block, left=of rl]         (out)  {Orchestra\\renderer};

  \draw[arrow] (in)    -- (clean);
  \draw[arrow] (clean) -- (hmm);
  \draw[arrow] (hmm)   -- (oltw);
  \draw[arrow] (oltw)  -- (rl);
  \draw[arrow] (rl)    -- (out);

  \draw[arrow, dashed] (hmm.south west) to[bend right=15]
    node[left, font=\tiny, pos=0.55]
        {$\max\alpha_t<\theta_{\text{conf}}$} (oltw.north west);

\end{tikzpicture}
}\par
\vspace{2pt}
\footnotesize
\figurename~\thefigure. Схематическое сопоставление двух
архитектур: классический пайплайн с входным модулем, очисткой,
HMM-трекером, OLTW-fall-back и рендерером~--- то, что
сохраняется в нашей постановке~\Secref{proposal}; и RL-замена
трекера в работе~\cite{Dorfer2018RL}, где RL-агент действует
на месте HMM/OLTW (штриховая стрелка).\par}
\addvspace{4pt} Вознаграждение строится через определение
истинной позиции $x$ исполнителя, рассчитанной через
ground-truth выравнивание; для текущей предсказанной позиции
$\hat x$ и сдвига $d_x = \hat x - x$ функция награды имеет
вид~\cite[\S~3]{Dorfer2018RL}
\begin{equation}
r_t = \begin{cases}
\,1 - |d_x|/b, & |d_x| \le b,\\
\;\text{сброс эпизода},   & |d_x| > b,
\end{cases}
\eqlab{dorfer-reward}
\end{equation}
где $b$~--- ширина окна допустимых отклонений. Дисконт $\gamma=0.9$.
Архитектура~--- общая для policy и value свёрточная сеть с
ELU-активациями и dropout~$0.2$, отдельные ветки обработки
аудио и изображения с конкатенацией признаков, выход policy~---
softmax по трём действиям, выход value~--- скалярный.
Обучение~--- A2C с шестнадцатью параллельными
акторами, $t_{\max}=15$~\cite[\S~4]{Dorfer2018RL}; сравнение с
REINFORCE-with-baseline проведено в той же работе.

С точки зрения позиционирования, эта схема представляет собой
\textit{RL-агента, заменяющего трекер целиком}: классические
HMM/OLTW-механизмы в систему не входят. Сильная сторона~---
концептуальная элегантность и end-to-end обучаемость; слабая~---
зависимость от тщательно размеченных пар изображение-онсет и
ограниченность экспериментов синтетическими данными
LilyPond+SoundFont. Авторы упоминают возможность расширения
действия на непрерывное и обобщения на полифонию как
направления дальнейшей работы.

\paragraph*{Хенкель и Видмер 2021: мультиразрешающее предсказание
по изображению.} Та же группа CPJKU развила
направление~\cite{Henkel2021}: задача формулируется уже не как
RL, а как \textit{детерминированная} задача регрессии
ограничивающих рамок (bounding boxes) объектов на изображении
партитуры с условием от аудио. Архитектура~--- conditional
YOLO с FiLM-conditioning, в котором аудио-вектор $z$ модулирует
карты признаков по правилу $f_{\text{FiLM}}(x) = s(z)\odot x +
t(z)$~\cite[уравн.~1]{Henkel2021}. Существенный методологический
вклад~--- одновременное обучение на разметках трёх уровней
(нота, такт, система) с весами потерь $\lambda^{(l)}_{\text{box}}$,
$\lambda^{(l)}_{\text{obj}}$. Хотя эта работа формально не
относится к RL, она представляет наиболее современный
end-to-end подход к отслеживанию партитуры по изображению и
является прямым продолжением~\cite{Dorfer2018RL} в supervised-режиме.

\paragraph*{Петер 2023: онлайн-выравнивание символьного входа
через offline RL.} Современная работа Петер~\cite{Peter2023}
предлагает альтернативный путь: классический оффлайновый
двухшаговый DTW-выравниватель решает задачу при доступности всей
последовательности, а онлайн-агент обучается \textit{предсказывать
тот же ответ} в каузальном режиме методами Q-learning. Состояние
агента~--- последовательность из 16 элементов: 7 элементов
прошлого партитуры, текущая оценка позиции, 8 элементов будущего
партитуры, плюс 8 последних нот исполнителя~\cite[\S~III]{Peter2023}.
Действие~--- выбор одного из 16 кандидатных onset-событий в окне.
Вознаграждение бинарное~--- 1 за правильный выбор, 0 за
неправильный. С дисконтом $\gamma=0$ (полностью myopic-агент)
квадратичная TD-функция потерь сводится к бинарной
кросс-энтропии: задача обучения становится формально
эквивалентной классификации правильности каждого кандидата.
Архитектура~--- трансформер на 8~голов, 6 слоёв,
$\sim$157 тысяч параметров. На корпусе симфонических
исполнений~\cite[\S~IV]{Peter2023}~автор сообщает медианную
асинхронность $15.7$\,мс для системы Online Alignment Model и
$36.0$\,мс для упрощённого Greedy Agent против $60.6$\,мс у
референсного OLTW; F1-оценка нот по разным произведениям
варьируется в пределах $90\,\%$--$96\,\%$.

С позиционной точки зрения эта работа предлагает гибридную
схему: существенно отличающуюся от чистой RL-замены трекера и
формально близкую к нашему предложению (см.~\Secref{proposal}),
однако сосредоточенную на \textit{символьном} входе MIDI и
полностью myopic-агенте. Прямой постановки опережающего
предсказания темпа (anticipation) с горизонтом $200$--$500$\,мс
в работе нет.

\paragraph*{Цзян и др. 2020: RL-Duet.} В работе
RL-Duet~\cite{Jiang2020RLDuet} рассматривается родственная,
но иная задача~--- генерация второго голоса контрапункта
к произведённой человеком монодической мелодии в режиме
реального времени. Состояние $s_t = (s_t^h, s_t^m)$ объединяет
скользящие окна последних токенов человеческой и машинной
партий; действие~--- следующий машинный токен из дискретного
алфавита. Вознаграждение собирается из ансамбля четырёх
обученных нейросетевых критиков: совместный pre/post-критик,
горизонтальный Cloze-критик на машинной партии, вертикальный
гармонический критик по человеческой партии, плюс слабое
наказание $-1$ за чрезмерные повторы (по
аналогии с~\cite{Jaques2017SeqTutor}). Алгоритм обучения~---
актор-критик с GAE~\cite{Schulman2016GAE}; дисконт
$\gamma=0.5$, $100\,000$~training duets. Существенно, что
постановка~\cite{Jiang2020RLDuet} \textit{не относится} к
отслеживанию партитуры~--- это задача \textit{генерации} нового
музыкального материала, не выравнивания с заранее известной
партитурой.

\paragraph*{Ву и др. 2024: ReaLchords и KL-контроль.} Развитием
линии RL-Duet является система ReaLchords~\cite{Wu2024ReaLchords},
в которой решается принципиальная проблема \textit{exposure bias}
в онлайн-генерации: при онлайн-сэмплировании каузальная политика
$\pi_\theta(y_t|x_{<t},y_{<t})$ накапливает ошибки, не
встречавшиеся в обучающих парах. Авторы предлагают двухстадийную
схему: сначала \textit{онлайн} MLE на парах $(x,y)$, затем
RL-дообучение целью~\cite[уравн.~2]{Wu2024ReaLchords}
\begin{equation}
\max_\theta \Expect_{x\sim p_{\text{data}},\,y\sim\pi_\theta}\!\bigl[
R(x,y) - \beta\,\mathrm{KL}(\pi_\theta\,\|\,\phi_\omega)\bigr],
\eqlab{realchords-obj}
\end{equation}
где $\phi_\omega(y|x)=\prod_t\phi_\omega(y_t|x,y_{<t})$~---
\textit{офлайновая} модель учителя, видящая всю мелодию~$x$
сразу, а $R(x,y)$~--- ансамбль контрастивных и дискриминативных
наград. Эта же работа подробно разбирает понятие
\textit{anticipation} применительно к гармонической генерации.
Постановка остаётся в области \textit{генерации} аккордов,
а не отслеживания партитуры. Однако формальное определение
двухстадийной схемы (онлайн-каузальная политика с офлайн-учителем)
методологически близко к гибридной схеме настоящей работы и
служит важным методологическим прецедентом.

\paragraph*{Жак и др. 2017: Sequence Tutor.} Sequence
Tutor~\cite{Jaques2017SeqTutor}~--- ранняя и общая работа по
\textit{KL-контролю} как способу дообучения предобученной
LSTM-модели последовательности с внешним эвристическим
вознаграждением. Авторы предлагают три алгоритмических варианта
($Q$-learning с log-prior augmentation, обобщённое $\Psi$-обучение
и $G$-learning) и применяют их к генерации монофонических MIDI-мелодий
с правилами вознаграждения, поощряющими тонические ноты, мотивные
повторения и наказывающими автокорреляцию. Работа представляет
интерес как методологический прецедент для применения KL-контроля
в музыкальных задачах и используется как baseline в более поздних
работах~\cite{Jiang2020RLDuet,Wu2024ReaLchords}; прямого отношения
к отслеживанию партитуры она не имеет.

\paragraph*{Сводное позиционирование.} Из проведённого обзора
видно, что существующие RL-подходы образуют четыре неперекрывающихся
группы: (i)~RL-агент \textit{заменяет} трекер целиком и обучается
end-to-end~\cite{Dorfer2018RL}; (ii)~RL-агент решает задачу
\textit{онлайн-выравнивания символьного} входа в myopic-режиме при
наличии офлайнового референса~\cite{Peter2023}; (iii)~RL применяется
к \textit{генерации} нового музыкального материала
(контрапункт, гармонизация)~\cite{Jiang2020RLDuet,Wu2024ReaLchords};
(iv)~RL применяется к общему дообучению последовательностной модели
с внешним вознаграждением~\cite{Jaques2017SeqTutor}. Ни одна из
указанных работ не рассматривает RL-агента как \textit{отдельный
тонкий слой опережающего предсказания темпа}, действующий поверх
сохраняемого классического вероятностного трекера. Именно этот пробел
обсуждается в~\Secref{proposal} как направление возможной новизны.
